% POYO+ terminal showcase — run sheet
% Build:  cd demo/guide && xelatex showcase_guide.tex && xelatex showcase_guide.tex
\documentclass[11pt]{article}

\usepackage[letterpaper, top=0.75in, bottom=0.8in, left=0.85in, right=0.85in, headheight=15pt]{geometry}
\usepackage{fontspec}
\setmainfont{Avenir Next}[BoldFont={* Demi Bold}, ItalicFont={* Italic}, BoldItalicFont={* Demi Bold Italic}]
\setsansfont{Avenir Next}[BoldFont={* Demi Bold}]
\setmonofont{FiraCode Nerd Font Mono}[Scale=MatchLowercase, Contextuals=NoAlternate]
\newfontfamily\symfont{Apple Symbols}[Scale=1.12]
\newfontfamily\headfont{Avenir Next}[UprightFont={* Demi Bold}, BoldFont={* Heavy}]

\usepackage{xcolor}
\usepackage[most]{tcolorbox}
\usepackage{graphicx}
\usepackage{booktabs}
\usepackage{tabularx}
\usepackage{array}
\usepackage{enumitem}
\usepackage{titlesec}
\usepackage{fancyhdr}
\usepackage{multicol}
\usepackage{microtype}
\usepackage{needspace}
\usepackage[hidelinks]{hyperref}
\usepackage{xurl}

% ---------------------------------------------------------------- numbers
% (filled in from the verification runs on this Mac)
\newcommand{\PretrainAvg}{0.84}
\newcommand{\PretrainBest}{0.91 (m\_20150610\_center\_out\_reaching)}
\newcommand{\MonkeyTAvg}{0.71}
\newcommand{\BCIRtwo}{0.76}
\newcommand{\DecodeSeconds}{25}
\newcommand{\DecodeWindows}{2{,}770}
\newcommand{\DecodeMinutes}{25}
\newcommand{\PretrainEpochs}{60}
\newcommand{\PretrainSteps}{7{,}200}
\newcommand{\PretrainHours}{1 h 44 min}
\newcommand{\PretrainSkipped}{73}
\newcommand{\LiveFinal}{0.27}
\newcommand{\LiveMinutes}{3.5}
\newcommand{\FinetuneMinutes}{2.5}
\newcommand{\FinetuneFrozen}{0.51}
\newcommand{\TestsPassed}{48}


% ---------------------------------------------------------------- colours
% Catppuccin Latte (on paper) + Catppuccin Macchiato (terminal boxes)
\definecolor{ltext}{HTML}{4C4F69}
\definecolor{lsub}{HTML}{6C6F85}
\definecolor{lmauve}{HTML}{8839EF}
\definecolor{lblue}{HTML}{1E66F5}
\definecolor{lgreen}{HTML}{40A02B}
\definecolor{lpeach}{HTML}{FE640B}
\definecolor{lred}{HTML}{D20F39}
\definecolor{lyellow}{HTML}{DF8E1D}
\definecolor{lteal}{HTML}{179299}
\definecolor{lsurface}{HTML}{CCD0DA}
\definecolor{lmantle}{HTML}{E6E9EF}
\definecolor{lbase}{HTML}{EFF1F5}
\definecolor{mbase}{HTML}{24273A}
\definecolor{mtext}{HTML}{CAD3F5}
\definecolor{msurface}{HTML}{494D64}
\definecolor{mmauve}{HTML}{C6A0F6}
\color{ltext}

% ---------------------------------------------------------------- headings
\titleformat{\section}{\headfont\LARGE\color{lmauve}}{\thesection}{0.6em}{}[\vspace{-0.6ex}{\color{lsurface}\titlerule[0.8pt]}]
\titleformat{\subsection}{\headfont\Large\color{lblue}}{\thesubsection}{0.5em}{}
\titleformat{\subsubsection}{\headfont\large\color{ltext}}{}{0em}{}
\titlespacing*{\section}{0pt}{2.2ex plus 1ex}{1.4ex}
\titlespacing*{\subsection}{0pt}{2ex plus 1ex}{0.9ex}
\setlength{\parindent}{0pt}
\setlength{\parskip}{0.55ex plus 0.2ex}
\setlist{itemsep=0.25ex, topsep=0.5ex}

\pagestyle{fancy}
\fancyhf{}
\fancyhead[L]{\small\color{lsub}POYO+ in the terminal · showcase run sheet}
\fancyhead[R]{\small\color{lsub}\thepage}
\renewcommand{\headrulewidth}{0pt}

% ---------------------------------------------------------------- keys & code
\newcommand{\key}[1]{\tcbox[on line, colback=lbase, colframe=lsub!70, boxrule=0.5pt, arc=2.2pt, left=2.6pt, right=2.6pt, top=0.7pt, bottom=0.7pt, boxsep=0pt]{\small\strut #1}}
\newcommand{\CMD}{{\symfont ⌘}}
\newcommand{\SHIFT}{{\symfont ⇧}}
\newcommand{\OPT}{{\symfont ⌥}}
\newcommand{\CTRL}{{\symfont ⌃}}
\newcommand{\RET}{{\symfont ⏎}}
\newcommand{\UP}{{\symfont ↑}}
\newcommand{\DOWN}{{\symfont ↓}}
\newcommand{\LEFT}{{\symfont ←}}
\newcommand{\RIGHT}{{\symfont →}}
\newcommand{\TAB}{{\symfont ⇥}}
\newcommand{\arr}{{\symfont →}}
\newfontfamily\dingfont{Zapf Dingbats}
\newcommand{\checkmark}{{\dingfont ✔}}
\newcommand{\Return}{\key{\RET\,return}}
\newcommand{\code}[1]{\texttt{\detokenize{#1}}}
\newcommand{\file}[1]{\textcolor{lteal}{\nolinkurl{#1}}}

\newtcblisting{term}[1][]{listing only, enhanced, colback=mbase, colframe=msurface, boxrule=0.6pt, arc=4pt,
  left=9pt, right=8pt, top=4pt, bottom=4pt, before skip=0.8ex, after skip=1.0ex,
  listing engine=listings,
  listing options={basicstyle=\ttfamily\small\color{mtext}, columns=fullflexible, keepspaces=true,
    breaklines=true, upquote=true, showstringspaces=false},
  overlay={\node[anchor=north east, font=\tiny\ttfamily, text=mmauve, inner sep=3pt] at (frame.north east) {fish};},
  #1}

\newtcolorbox{say}{enhanced, colback=lmauve!5, colframe=lmauve!60, boxrule=0pt, leftrule=3pt, arc=0pt,
  left=8pt, right=6pt, top=4pt, bottom=4pt, before skip=1ex, after skip=1ex,
  fontupper=\itshape\color{ltext}, title={\upshape\bfseries\small SAY}, attach title to upper={\ \ },
  coltitle=lmauve}
\newtcolorbox{tip}[1][Tip]{enhanced, colback=lgreen!6, colframe=lgreen!70, boxrule=0pt, leftrule=3pt, arc=0pt,
  left=8pt, right=6pt, top=4pt, bottom=4pt, before skip=1ex, after skip=1ex,
  title={\bfseries\small #1}, attach title to upper={\ \ }, coltitle=lgreen}
\newtcolorbox{warn}[1][Heads-up]{enhanced, colback=lpeach!6, colframe=lpeach!80, boxrule=0pt, leftrule=3pt, arc=0pt,
  left=8pt, right=6pt, top=4pt, bottom=4pt, before skip=1ex, after skip=1ex,
  title={\bfseries\small #1}, attach title to upper={\ \ }, coltitle=lpeach}
\newtcolorbox{sees}{enhanced, colback=lblue!4, colframe=lblue!60, boxrule=0pt, leftrule=3pt, arc=0pt,
  left=8pt, right=6pt, top=4pt, bottom=4pt, before skip=1ex, after skip=1ex,
  title={\bfseries\small WHAT THE AUDIENCE SEES}, attach title to upper={\\[0.3ex]}, coltitle=lblue}

% act header: number, title, duration
\newcommand{\act}[3]{%
  \clearpage
  \begin{tcolorbox}[enhanced, colback=mbase, colframe=mbase, arc=5pt, left=12pt, right=12pt, top=8pt, bottom=8pt]
    {\headfont\color{mmauve}\large ACT #1}\hfill{\small\color{mtext}#3}\\[0.3ex]
    {\headfont\color{mtext}\huge #2}
  \end{tcolorbox}
  \addcontentsline{toc}{section}{Act #1 · #2}%
}

\newcommand{\shot}[2][0.95]{%
  \IfFileExists{#2}{\begin{center}\includegraphics[width=#1\linewidth, height=0.80\textheight, keepaspectratio]{#2}\end{center}}{%
  \begin{center}\color{lred}[missing image: \detokenize{#2}]\end{center}}}

\begin{document}

% ================================================================= cover
\thispagestyle{empty}
\begin{tcolorbox}[enhanced, colback=mbase, colframe=mbase, arc=6pt, left=18pt, right=18pt, top=16pt, bottom=16pt]
  {\headfont\color{mmauve}\small TORCH\_BRAIN · TERMINAL SHOWCASE}\\[1.2ex]
  {\headfont\color{mtext}\fontsize{34}{38}\selectfont POYO+ in the terminal}\\[1.4ex]
  {\color{mtext}\Large A keystroke-by-keystroke run sheet for demoing a transformer that decodes
  monkey motor-cortex spikes into hand movements, live on your M1 Max, in your iTerm2.}
\end{tcolorbox}

\vspace{1ex}
\shot[1.0]{img/act1_banner.pdf}

\begin{tcolorbox}[enhanced, colback=lbase, colframe=lsurface, arc=4pt, left=10pt, right=10pt]
\textbf{The 10-second version.} Open iTerm2 \arr{} \code{z torch} \arr{} \code{conda activate torch_brain}
\arr{} \code{zsh demo/showcase.sh} and press \Return{} between acts. Everything below explains each act in
detail, what to type, what to say, and which terminal tricks to show off along the way.
\end{tcolorbox}

\begin{center}
\renewcommand{\arraystretch}{1.25}
\begin{tabularx}{\linewidth}{>{\bfseries\color{lmauve}}l X >{\color{lsub}}r}
\toprule
\textcolor{ltext}{Act} & \textbf{What happens} & \textcolor{ltext}{\textbf{Time}} \\
\midrule
0 & Set the stage: fastfetch, zoxide, starship prompt, conda, eza, git, glow & 3 min \\
1 & Meet the data: 111 recordings, a truecolour spike heatmap, spikes \arr{} tokens & 1 min \\
2 & Inside POYO+: architecture tree, a real forward pass, speed benchmark & 2 min \\
3 & Train POYO+ live from scratch (split pane + GPU monitor) & \LiveMinutes{} min \\
4 & Decode every held-out trial with the pretrained model (R² \PretrainAvg) & 1 min \\
5 & Live brain-computer-interface replay, drawn in braille & 1 min \\
6 & Transfer to monkey T, an animal the model has never seen & \FinetuneMinutes{} min \\
7 & Decode monkey T + live BCI on monkey T (R² \MonkeyTAvg) & 2 min \\
8 & Encore: tests, TensorBoard, lazygit, nvim, iTerm2 Instant Replay & 2--5 min \\
\bottomrule
\end{tabularx}
\end{center}
{\small\color{lsub} Full show $\approx$ 20--25 minutes. Short show (skip Acts 3 and 6, which train models live): $\approx$ 10 minutes.}

\clearpage
\tableofcontents

% ================================================================= what is included
\section{What's in this demo}

Everything needed for the demo is installed, trained and tested on this Mac
(Apple M1 Max, 64 GB, macOS). Nothing needs to be downloaded on show day.

\subsection*{Software}
\begin{itemize}
\item \textbf{Python environment:} the existing conda env \code{torch_brain} (Python 3.11) now contains
  \code{torch_brain} (installed in editable mode, \code{pip install -e ".[dev]"}) with the versions pinned in
  \file{setup.py}: PyTorch 2.2.0 (with Apple-GPU/MPS support), Lightning 2.3.3, Hydra 1.3, temporaldata 0.1.1, plus
  the demo extras in \file{demo/requirements.txt} (torchmetrics 1.6.1, NumPy 1.23, TensorBoard, plotext 5.3.2, \dots).
\item \textbf{Homebrew tools added for the show:} \code{fastfetch} (system splash), \code{btop} (resource monitor),
  \code{macmon} (Apple-Silicon CPU/GPU/power monitor, no sudo needed), \code{glow} (Markdown in the terminal).
  Already present and used: \code{bat}, \code{eza}, \code{fzf}, \code{zoxide}, \code{starship}, \code{lazygit}, \code{nvim}.
\item \textbf{Your shell:} the Research iTerm2 profile starts \textbf{fish} with your Terminal-Coding-Configs setup
  (starship prompt, zoxide, fzf, eza abbreviations), not zsh. Every command in this guide was tested in that shell.
  zoxide already knows the repo, so \code{z torch} jumps to it from anywhere.
\end{itemize}

\subsection*{Data and models}
\begin{itemize}
\item \textbf{Data:} your \file{data/processed/perich_miller_population_2018} (111 sessions, 2.84 GB) was already processed
  by brainsets and is used as-is.
\item \textbf{Pretrained model:} \file{demo/checkpoints/poyo_plus_demo.ckpt}, a 5.2 M-parameter POYO+ trained for
  \PretrainEpochs{} epochs (\PretrainSteps{} steps, \PretrainHours) on the Apple GPU, on 12 recordings from monkeys C, M and J.
  Average held-out R² = \textbf{\PretrainAvg}.
\item \textbf{Monkey-T model:} \file{demo/checkpoints/poyo_plus_monkey_t.ckpt}, the pretrained model transferred to a
  session of monkey T (never seen during pretraining). Held-out R² = \textbf{\MonkeyTAvg}.
  Act 6 re-creates it live; this copy is the safety net.
\end{itemize}

\subsection*{Code fixes and improvements}
\renewcommand{\arraystretch}{1.2}
\begin{tabularx}{\linewidth}{>{\raggedright\arraybackslash}p{4.3cm} >{\raggedright\arraybackslash}X}
\toprule
\textbf{Where} & \textbf{What changed and why} \\
\midrule
\file{train.py} & Default config pointed to a non-existent \code{train.yaml} (now \code{train_poyo_mp}); device is chosen
  automatically (CUDA \arr{} Apple GPU \arr{} CPU); TensorBoard logging when W\&B is off; optional
  Rich progress bars; gradient clipping + skip of non-finite steps; iTerm2 notification when training ends. \\
\file{torch_brain/models/poyo_plus.py} & Timestamps are emitted as float32: Apple GPUs (MPS) cannot hold float64 tensors, so
  training on the GPU crashed. \\
\file{torch_brain/nn/embedding.py}, \file{infinite_vocab_embedding.py} & MPS-only, mathematically identical embedding
  paths. Padding tokens all scattered their gradients into one row and serialised on the GPU (0.6 s of a 1 s step).
  Training steps are now $\sim$1.7$\times$ faster. \\
\file{torch_brain/optim.py} & Sync-free \code{Lamb} (same maths as \code{torch_optimizer.Lamb}, verified to 1e-8):
  removes 2--3 GPU$\to$CPU syncs per parameter per step. \\
\file{torch_brain/utils/stitcher.py} & \textbf{Bug fix:} every window of a validation batch was credited to the
  \emph{first} sequence of the batch, mixing recordings in the per-session R². Also: CPU bookkeeping (much faster
  validation), float64 time stitching, Rich R² table. \\
\file{torch_brain/nn/multitask_readout.py} & New \code{ignore_intervals} option: annotated cursor-tracking glitches
  (up to 1300 cm/s!) are ignored by loss and metrics. This doubled R² for the same compute and removed NaN blow-ups.
  Faster \code{unpack_output}. \\
\file{finetune.py} (new, root) & Working version of the broken \file{examples/poyo_plus/finetune.py}
  (imported a non-existent \code{StitchEvaluator}). \\
\file{tests/} & Two new test files (8 tests); the full suite passes (\TestsPassed{} passed). \\
\bottomrule
\end{tabularx}

\subsection*{New files}
\begin{tabularx}{\linewidth}{>{\raggedright\arraybackslash}p{5.2cm} >{\raggedright\arraybackslash}X}
\toprule
\file{demo/01_explore_data.py} & Act 1 · data tour \\
\file{demo/02_model_tour.py} & Act 2 · architecture + forward pass + benchmark \\
\file{demo/03_decode.py} & Acts 4/7 · decode held-out trials, R² table, braille chart, figure \\
\file{demo/04_live_bci.py} & Acts 5/7 · live BCI replay \\
\file{demo/showcase.sh} & autopilot: runs every act, waits for \key{\RET} \\
\file{demo/check_setup.py} & preflight checklist \\
\file{demo/README.md}, \file{demo/setup_env.sh} & overview (rendered with glow) · env re-creation \\
\file{configs/demo_*.yaml}, \file{configs/model/poyo_demo.yaml}, \file{configs/dataset/demo_*.yaml} & Hydra configs for the demo \\
\file{torch_brain/utils/terminal.py} & iTerm2 inline images, notifications, fireworks, badges, tab colours \\
\bottomrule
\end{tabularx}

% ================================================================= before the show
\section{Before the show (T--15 minutes)}

\subsection{Prepare the machine}
\begin{enumerate}[leftmargin=*]
\item \textbf{Plug in the charger.} Training is GPU-heavy; on battery macOS may throttle it.
\item \textbf{Quit GPU-hungry apps} (games, video editors, many Chrome tabs, Docker). They slow the live training acts.
\item \textbf{Notifications.} iTerm2 is set to post a macOS notification when training finishes. Make sure
  \emph{System Settings \arr{} Notifications \arr{} iTerm2 \arr{} Allow Notifications}
  is on, and that a Focus/Do-Not-Disturb mode is \emph{off} (or the banner will be silenced).
  The fireworks and the tab colours work regardless.
\end{enumerate}

\subsection{Open and size iTerm2}
\begin{enumerate}[leftmargin=*]
\item Press \key{\CMD\,space}, type \code{iTerm}, press \Return. Your default profile is \textbf{Research}
  (fish, Iosevka Nerd Font 14 pt, Catppuccin Macchiato colours); all demo output is themed to match it.
\item Press \key{\CMD\,\RET} to go full screen (press it again to leave).
\item Adjust the font for the room with \key{\CMD\,+} / \key{\CMD\,--} (\key{\CMD\,0} resets). The live BCI act
  needs at least $\sim$\textbf{115 columns $\times$ 42 rows}. Check with:
\end{enumerate}
\begin{term}
tput cols; tput lines
\end{term}

\needspace{9\baselineskip}
\subsection{Run the preflight check}
Type each line and press \Return{} after it:
\begin{term}
z torch
conda activate torch_brain
python demo/check_setup.py
\end{term}
Every \emph{required} row must show a green \textbf{\checkmark}. Optional rows (yellow \textbf{!}) only disable eye candy.
The iTerm2 row must be green when you run it inside iTerm2.

\begin{tip}[Rehearse once]
Before the audience arrives, type each act's command yourself once (or at least the acts you plan to show). That warms
the GPU and file caches, and it puts every command into fish's history, so during the show fish offers the rest
of each command in grey after a couple of letters. Press \key{\RIGHT} to accept it.
That looks slick and avoids typos. \key{\CMD\,R} (or \key{\CTRL\,R}) opens fzf history search to find them instantly.
(Commands run by the autopilot script do not go into your history.)
\end{tip}

Finally type \code{cd} \Return{} to go back to your home folder (fish's \code{cd} with no argument goes home), then
press \key{\CMD\,K} to clear the screen. You are ready.

\begin{warn}[Start Act 0 from outside the repo]
\code{z torch} skips the folder you are already in. Typed inside the repo root, it jumps to the next match instead: the
package folder \file{torch_brain/torch_brain}, where \code{python demo/...} commands fail with ``No such file''. That's
why you go home first. If it happens anyway, type \code{cd ..} \Return.
\end{warn}

% ================================================================= act 0
\act{0}{Set the stage}{3 min}

\begin{say}
This is my terminal: iTerm2 running fish, with a starship prompt, zoxide, fzf and a few modern CLI tools. Everything you will
see today, including training a neural network on monkey brain recordings, runs right here.
\end{say}

\begin{enumerate}[leftmargin=*]
\item \textbf{System splash.} Type \code{fastfetch} \Return. It shows the Apple logo, macOS, the M1 Max chip and GPU,
  memory, the shell, terminal and font.
\item \textbf{Jump to the project with zoxide.} Type \code{z torch} \Return. The blue directory segment of the prompt
  changes to \code{~/Documents/GitHub/torch_brain}, and a mauve git segment appears with the current branch
  (\code{poyo-terminal-demo}) and file counts such as \code{new:1}. The teal segment shows the Research project,
  and the right side shows a green \textbf{\checkmark} for the last command's exit status.
\item \textbf{Activate the environment.} Type \code{conda activate torch_brain} \Return. The prompt gains a
  green segment that names the environment.
\begin{tip}[fish syntax highlighting]
While typing, watch the command word: it is red while it isn't a valid command and turns blue once it is. Type
\code{cond} slowly (red), then finish \code{conda} (blue).
\end{tip}
\item \textbf{Look around with eza.} Your Research config has abbreviations that expand as you type:
\begin{term}
ll
tree --git-ignore demo
\end{term}
  Type \code{ll} and press \key{space}: fish expands it in place to the full \code{eza} command (long listing with a
  header, a git-status column, directories first). Press \Return{} to run it. \code{tree} expands the same way to a two-level eza tree;
  \code{--git-ignore} hides cache folders.
\item \textbf{The work behind this demo, in git:}
\begin{term}
git log --oneline -4
\end{term}
  Three commits on top of the original code: the Apple-GPU and evaluation fixes, training scripts at the repo root, and
  this demo with its run sheet.
\item \textbf{Read the project overview, rendered in the terminal:}
\begin{term}
glow -p demo/README.md
\end{term}
  It scrolls in a pager: \key{space} next page, \key{b} back, \key{q} quit. (Drop \code{-p} to print it all at once.)
\end{enumerate}

\begin{say}
POYO+ is a transformer for neuroscience. It treats every single spike as a token, so recordings with different
neurons, or even different animals, can all feed the same network.
\end{say}

% ================================================================= act 1
\act{1}{Meet the data}{1 min}

\begin{term}
python demo/01_explore_data.py
\end{term}
{\small\color{lsub}Runtime $\approx$ 6 s. Add \code{--fast} to skip the scan of all 111 files.}

\begin{sees}
\begin{enumerate}[leftmargin=*, nosep]
\item The \textbf{POYO+ banner} in a Catppuccin gradient. iTerm2 also shows a translucent \textbf{badge} (``POYO+ / ACT 1'')
  in the corner of the pane.
\item A spinner and progress bar while all \textbf{111 recordings} are scanned, then a table: 4 monkeys,
  111.4 M spikes, 43 hours of recordings, 2.84 GB on disk. Monkey T is marked \emph{held out}.
\item The dataset description and its DANDI link (\key{\CMD}-click it to open the archive in the browser).
\item \textbf{Live from motor cortex:} a truecolour heatmap of the 22 most active neurons over 10 seconds, drawn
  entirely with coloured text characters.
\item \textbf{Spikes \arr{} tokens:} the first spike tokens of a 1-second window (neuron id, time), and the token budget:
  $\sim$600 spike tokens \arr{} 128 latents \arr{} 100 velocity outputs.
\item A high-resolution figure \textbf{inline in the terminal} (iTerm2 inline images): spike raster, cursor
  velocity and the beautiful 8-direction ``flower'' of every reach in the session.
\end{enumerate}
\end{sees}

\begin{say}
Four monkeys were trained to move a cursor to targets while the Miller lab recorded from up to a few hundred neurons in motor
cortex. The goal is to read those spikes and predict the hand's velocity in real time.
\end{say}


\begin{tip}[iTerm2 marks]
Every act drops an iTerm2 \emph{mark}. Later, \key{\CMD\,\SHIFT\,\UP} / \key{\CMD\,\SHIFT\,\DOWN} jump back and forth between
them, which is handy to scroll back to the heatmap or a figure while answering questions.
\end{tip}

\shot[0.55]{img/act1.pdf}
\shot[1.0]{img/fig_01_data_overview.png}

% ================================================================= act 2
\act{2}{Inside POYO+}{2 min}

First show the heart of the model's code, syntax-highlighted in your theme:
\begin{term}
bat --theme="Catppuccin Macchiato" -r 199:245 torch_brain/models/poyo_plus.py
\end{term}
{\small\color{lsub}These 47 lines are the whole forward pass: embed spikes \arr{} \code{# encode} (cross-attention)
\arr{} \code{# process} (self-attention) \arr{} \code{# decode} \arr{} readout. If bat opens a pager, press \key{q}.}

Then run the tour:
\begin{term}
python demo/02_model_tour.py
\end{term}
{\small\color{lsub}Runtime $\approx$ 10 s.}

\begin{sees}
\begin{enumerate}[leftmargin=*, nosep]
\item \textbf{Architecture tree} with parameter counts: input tokens (unit embedding for 801 neurons, token types,
  rotary time), encoder, 6-layer processor, decoder (session + task embeddings), multitask readout.
\item \textbf{A real batch through the network:} 8 random one-second windows from the 12 recordings, and the tensor
  shape at each stage (padded spike tokens \arr{} 128 latents \arr{} 100 queries \arr{} (v\textsubscript{x}, v\textsubscript{y})).
\item \textbf{Speed:} 8 seconds of brain activity decoded in a few milliseconds on the GPU, and on the CPU for
  comparison: hundreds of times faster than real time.
\end{enumerate}
\end{sees}

\begin{say}
No binning, no fixed channel order. Each neuron gets a learned embedding, and a fixed set of latent tokens
summarises the population every 125 ms. That's why the same model works across sessions and animals.
\end{say}

\shot[0.7]{img/act2.pdf}

% ================================================================= act 3
\act{3}{Train POYO+ live, from scratch}{\LiveMinutes{} min}

This act trains a fresh POYO+ on two recordings while the audience watches the GPU work.

\begin{enumerate}[leftmargin=*]
\item \textbf{Split the window:} press \key{\CMD\,D}. A second pane opens on the right.
\item \textbf{GPU monitor in the right pane:} type \code{macmon} \Return. It shows live CPU/GPU utilisation,
  frequency, power and RAM of the M1 Max. (Alternative: \code{btop}.)
\item \textbf{Back to the left pane:} press \key{\CMD\,[} (or \key{\CMD\,\OPT\,\LEFT}).
\item \textbf{Start training:}
\begin{term}
python train.py --config-name demo_live
\end{term}
\end{enumerate}

\begin{sees}
\begin{enumerate}[leftmargin=*, nosep]
\item The iTerm2 \textbf{tab turns yellow} (training in progress) and a Rich \textbf{model summary} table appears.
\item A colourful \textbf{progress bar} per epoch with the running \code{train_loss}. In the right pane the
  \textbf{GPU jumps to $\sim$90--100\%}.
\item After every epoch an \textbf{R² table} for the two held-out recordings, green/yellow/red bars that grow from
  $\approx$0 to $\approx$\LiveFinal{} after 8 epochs.
\item At the end the tab turns \textbf{green}, a macOS notification ``POYO+ finished training'' appears, and
  iTerm2 plays \textbf{fireworks} at the cursor. The right side of the prompt shows how long the command took
(e.g.\ \code{took 3m29s}) next to a green \textbf{\checkmark} (shown when the window is at least 120 columns wide).
\end{enumerate}
\end{sees}

\begin{say}
Every spike in a one-second window is a token. The loss is the error on hand velocity. Watch R² climb: it's
learning which neurons care about which direction, from scratch, in a few minutes on a laptop.
\end{say}

\begin{tip}[Hydra overrides: change anything from the command line]
Shorter run: \code{python train.py --config-name demo_live epochs=4}. Smaller model:
\code{model.depth=2}. CPU instead of GPU: \code{accelerator=cpu} (great for showing why the GPU matters).
\end{tip}

When it's done: press \key{\CMD\,]} to go to the right pane, \key{q} to quit macmon, then type \code{exit} \Return{}
to close the pane.

\begin{warn}
Don't press \key{\CTRL\,C} during the first seconds of training unless you mean it. It stops the run cleanly
(Lightning saves a checkpoint), but the act is over. The pretrained model used next does not depend on this run.
\end{warn}

% ================================================================= act 4
\act{4}{Decode held-out trials}{1 min}

\begin{term}
python demo/03_decode.py
\end{term}
{\small\color{lsub}Runtime $\approx$ \DecodeSeconds{} s. Uses \file{demo/checkpoints/poyo_plus_demo.ckpt}.}

\begin{sees}
\begin{enumerate}[leftmargin=*, nosep]
\item A progress bar while every held-out test trial of all 12 recordings is cut into overlapping 1-second windows
  and decoded (\DecodeWindows{} windows, \DecodeMinutes{} minutes of neural data).
\item \textbf{The R² leaderboard:} one row per recording (monkeys C, M, J; center-out and random-target tasks), sorted,
  with coloured bars. \textbf{Average R² = \PretrainAvg}; best \PretrainBest.
\item \textbf{A chart drawn in braille characters right in the terminal}: true vs decoded horizontal velocity
  for four back-to-back test trials.
\item \textbf{An inline figure:} velocity traces (green = real, purple = POYO+) and, on the right, the reach
  trajectories rebuilt by integrating the decoded velocity: the ``flower'' again, now \emph{from brain activity alone}.
\end{enumerate}
\end{sees}

\begin{say}
These trials were never seen in training. Overlapping windows are stitched back together and scored with R²,
where 1.0 is perfect and 0 is just guessing the average. On the right, solid lines are the reaches the model
reconstructs from spikes; dashed lines are what the monkey actually did.
\end{say}

\shot[0.78]{img/act4.pdf}
\shot[0.9]{img/fig_03_decoding.png}

% ================================================================= act 5
\act{5}{Live brain-computer interface}{1 min}

\begin{term}
python demo/04_live_bci.py
\end{term}
{\small\color{lsub}Runtime $\approx$ 40 s (8 trials in real time). Options: \code{--trials 4}, \code{--speed 0.5}
(slow motion), \code{--session c_20150701}.}

\begin{sees}
Held-out reaches are replayed \textbf{in real time}. Every 50 ms POYO+ reads only the \emph{previous} second of spikes
(causal decoding, like a real BCI) and moves the \textbf{purple} cursor; the \textbf{green} cursor is the monkey's real
hand. The canvas is drawn with braille dots (2$\times$4 pixels per character); the target lights up in its colour.
On the right: trial, target direction, phase (hold/REACH), population firing rate, true vs decoded speed and a
running R², plus a scrolling activity heatmap of the 24 most active neurons. At the end: the causal R² over all
replayed trials (\BCIRtwo).
\end{sees}

\begin{say}
This is what a brain-computer interface does: no hand needed. The purple cursor is driven purely by neurons firing,
decoded by a transformer, 20 times per second.
\end{say}

\begin{tip}[iTerm2 Instant Replay]
After the animation finishes, press \key{\CMD\,\OPT\,B}. iTerm2's \emph{Instant Replay} lets you scrub back in time
with \key{\LEFT}/\key{\RIGHT} and replay the animation frame by frame. Press \key{esc} to return.
\end{tip}

\shot[0.85]{img/act5.pdf}

% ================================================================= act 6
\act{6}{Transfer to a never-seen monkey}{\FinetuneMinutes{} min}

\begin{term}
python finetune.py --config-name demo_finetune
\end{term}

\begin{sees}
\begin{enumerate}[leftmargin=*, nosep]
\item The pretrained weights are loaded and \textbf{55 brand-new neurons} from monkey T get fresh embeddings.
\item Epochs 0--5: log line ``POYO+ Perceiver frozen'': only the new neuron and session embeddings learn
  (\emph{unit identification}). R² climbs from 0 to $\approx$\FinetuneFrozen{} although the transformer itself
  is untouched: the pretrained network already ``speaks motor cortex''; it only has to learn who is who.
\item Epoch 6: ``Perceiver unfrozen'': everything is fine-tuned and R² climbs further, to $\approx$\MonkeyTAvg{}
  after 16 epochs (a small dip right after unfreezing is normal).
\item The final model is copied to \file{demo/checkpoints/poyo_plus_monkey_t.ckpt}; tab turns green, notification,
  fireworks.
\end{enumerate}
\end{sees}

\begin{say}
Monkey T was never seen during pretraining: different animal, different electrodes, different neurons. We freeze the
network and only learn ``who is who'' among the new neurons. A few minutes later we can decode this new brain.
That's the promise of foundation models for neuroscience.
\end{say}

\begin{tip}[Short on time?]
Skip this act. The monkey-T checkpoint already exists, so Act 7 works without it.
\end{tip}

% ================================================================= act 7
\act{7}{Decode monkey T}{2 min}

\begin{term}
python demo/03_decode.py --monkey-t
python demo/04_live_bci.py --monkey-t
\end{term}

\begin{sees}
The same R² table and figures as Act 4, now for monkey T's held-out trials (\textbf{R² = \MonkeyTAvg}), and the
live BCI replay on the new animal.
\end{sees}

\shot[0.9]{img/fig_05_monkey_t_decoding.png}
\shot[0.8]{img/act7_bci.pdf}

% ================================================================= act 8
\act{8}{Encore}{2--5 min}

Pick whatever fits your audience.

\subsubsection*{The test suite}
\begin{term}
python -m pytest tests/test_poyo_plus.py tests/test_apple_silicon.py -v -W ignore -p no:cacheprovider
\end{term}
Green \code{PASSED} lines, including the new tests proving the Apple-GPU code paths are mathematically identical to
the originals. (Whole suite: \code{python -m pytest -q -W ignore}.)

\subsubsection*{TensorBoard dashboard}
\begin{term}
tensorboard --logdir demo/runs
\end{term}
\key{\CMD}-click the printed \code{http://localhost:6006/} link. It shows loss curves, learning rate, R² per recording,
gradient norms, and the pretraining run plus Acts 3 and 6. Stop the server with \key{\CTRL\,C}.

\subsubsection*{Browse the code changes}
\begin{itemize}[leftmargin=*]
\item \code{lg} (your abbreviation for \code{lazygit}): \key{\UP}/\key{\DOWN} select a file, its diff appears on the right; \key{q} quits.
\item \code{nvim torch_brain/models/poyo_plus.py}: search with \code{/class POYOPlus} \Return, quit with \code{:q} \Return.
\item \code{bat --diff torch_brain/utils/stitcher.py}: only the lines changed for the demo (including the stitching
  bug fix), with git markers in the gutter.
\end{itemize}

\subsubsection*{iTerm2 and fish tricks worth showing}
\renewcommand{\arraystretch}{1.25}
\begin{tabularx}{\linewidth}{>{\raggedright\arraybackslash}p{4.4cm} >{\raggedright\arraybackslash}X}
\toprule
\key{\CMD\,\SHIFT\,\UP} / \key{\CMD\,\SHIFT\,\DOWN} & jump between command marks (every act sets one) \\
\key{\CMD\,\OPT\,B} & Instant Replay: scrub back in time through the terminal's output \\
\key{\CMD\,\SHIFT\,H} & paste history \\
\key{\CMD\,;} & autocomplete from words already on screen \\
\key{\CMD\,F} / \key{\CMD\,/} & find in output / find the cursor \\
\key{\CMD\,D}, \key{\CMD\,\SHIFT\,D}, \key{\CMD\,\SHIFT\,\RET} & split vertically, split horizontally, maximise the current pane \\
\key{\CMD}-click & open a URL printed in the terminal (e.g.\ TensorBoard's link) \\
\key{\CMD\,R} or \key{\CTRL\,R} & fzf fuzzy history search (type \code{decode}, \key{\RET}) \\
\key{\CMD\,P} or \key{\CTRL\,T} & fzf file picker with a bat preview; e.g.\ type \code{bat} and a space, press \key{\CMD\,P}, pick a file, \key{\RET} \\
\key{\RIGHT} & accept fish's grey autosuggestion \\
\code{z torch}, \code{zi} & zoxide: jump to a remembered directory / pick one interactively with fzf \\
\code{ll}, \code{la}, \code{tree}, \code{lg} & abbreviations that expand in place when you press \key{space} \\
\bottomrule
\end{tabularx}

% ================================================================= autopilot
\section{Autopilot mode}

The script \file{demo/showcase.sh} runs every act in order. Before each act it prints the act's title and the exact
command, then waits for \key{\RET}.
\begin{term}
zsh demo/showcase.sh
zsh demo/showcase.sh 4
SKIP_TRAIN=1 zsh demo/showcase.sh
\end{term}
The first runs everything; the second starts at Act 4; the third skips the two acts that train models live
(3 and 6). The script is written for zsh, which is why each line starts with \code{zsh}; type them in your fish
shell exactly as shown. Press \key{\CTRL\,C} at any prompt to stop.

% ================================================================= troubleshooting
\clearpage
\section{Troubleshooting}
\renewcommand{\arraystretch}{1.3}
\begin{tabularx}{\linewidth}{>{\raggedright\arraybackslash}p{4.6cm} >{\raggedright\arraybackslash}X}
\toprule
\textbf{Symptom} & \textbf{Fix} \\
\midrule
\code{can't open file} with a path containing \file{torch_brain/torch_brain/demo} & You are in the package folder, one level too deep (\code{z torch} typed from inside the repo root skips the root). Type \code{cd ..} \Return. \\
\code{z torch} goes to some other folder & Run \code{cd ~/Documents/GitHub/torch_brain} once; zoxide learns every directory you \code{cd} into and ranks the ones you use most. \\
\code{fish: Unknown command} for a zsh habit (\code{j}, \code{gst}, \dots) & The Research shell is fish with its own config: use \code{z} for jumping, \code{git status} for git, \code{ll} for listings. \\
\code{ModuleNotFoundError} or wrong Python & The env isn't active: \code{conda activate torch_brain}. Check with \code{which python}. \\
Figures print ``saved to \dots'' instead of appearing & You are not in iTerm2 (e.g. VS Code's terminal or plain tmux). Open the file with the printed \code{open ...} command, or run in iTerm2. \\
The live BCI frame flickers or wraps & Window too small: go full screen (\key{\CMD\,\RET}), reduce the font (\key{\CMD\,--}), or maximise the pane (\key{\CMD\,\SHIFT\,\RET}). \\
Training is much slower than expected & Another job is using the GPU (check macmon). Kill leftovers with \code{pkill -f train.py}. Plug in the charger. \\
No notification at the end of training & macOS notification permission for iTerm2 or a Focus mode (see ``Before the show''). Fireworks still play. \\
\code{FileNotFoundError} for \file{poyo_plus_monkey_t.ckpt} & Run Act 6 once: \code{python finetune.py --config-name demo_finetune}. \\
Pretrained checkpoint lost & Re-train it (\PretrainHours{} on the M1 Max): \code{python train.py --config-name demo_pretrain}, then copy \file{demo/runs/poyo_plus_demo_pretrain/version_*/checkpoints/last.ckpt} to \file{demo/checkpoints/poyo_plus_demo.ckpt}. \\
Environment broken & Re-create it: \code{zsh demo/setup_env.sh}. \\
Tab stays green / badge stays on after the show & Run the reset command shown below the table (the autopilot does it automatically at the end). \\
Anything else & \key{\CTRL\,C} stops any act safely. Re-run the preflight: \code{python demo/check_setup.py}. \\
\bottomrule
\end{tabularx}

Reset the iTerm2 tab colour and badge:
\begin{term}
printf '\e]6;1;bg;*;default\a\e]1337;SetBadgeFormat=\a'
\end{term}

% ================================================================= appendix
\clearpage
\section{Appendix: the numbers behind the demo}

\subsection*{Pretrained model: held-out R² per recording}
\begin{center}\renewcommand{\arraystretch}{1.12}
\begin{tabular}{l l l r r l}\toprule
\textbf{Subject} & \textbf{Date} & \textbf{Task} & \textbf{Neurons} & \textbf{R²} & \\ \midrule
Monkey M & 2015-06-10 & center-out & 49 & 0.910 & \textcolor{lmauve}{\rule{3.64cm}{1.6ex}} \\
Monkey M & 2015-06-16 & center-out & 54 & 0.875 & \textcolor{lmauve}{\rule{3.50cm}{1.6ex}} \\
Monkey C & 2015-07-01 & center-out & 49 & 0.862 & \textcolor{lmauve}{\rule{3.45cm}{1.6ex}} \\
Monkey C & 2015-06-30 & center-out & 44 & 0.857 & \textcolor{lmauve}{\rule{3.43cm}{1.6ex}} \\
Monkey M & 2014-01-14 & random-target & 165 & 0.851 & \textcolor{lmauve}{\rule{3.41cm}{1.6ex}} \\
Monkey M & 2015-06-11 & center-out & 55 & 0.846 & \textcolor{lmauve}{\rule{3.38cm}{1.6ex}} \\
Monkey C & 2015-11-09 & center-out & 65 & 0.845 & \textcolor{lmauve}{\rule{3.38cm}{1.6ex}} \\
Monkey C & 2013-10-03 & center-out & 71 & 0.831 & \textcolor{lmauve}{\rule{3.32cm}{1.6ex}} \\
Monkey C & 2015-06-29 & center-out & 49 & 0.830 & \textcolor{lmauve}{\rule{3.32cm}{1.6ex}} \\
Monkey C & 2013-10-09 & random-target & 74 & 0.793 & \textcolor{lmauve}{\rule{3.17cm}{1.6ex}} \\
Monkey J & 2016-04-05 & center-out & 38 & 0.768 & \textcolor{lmauve}{\rule{3.07cm}{1.6ex}} \\
Monkey C & 2013-10-11 & random-target & 88 & 0.764 & \textcolor{lmauve}{\rule{3.06cm}{1.6ex}} \\
\midrule \textbf{Average} & & & & \textbf{0.836} & \\ \bottomrule
\end{tabular}\end{center}


\subsection*{Training recipe}
\begin{itemize}
\item \textbf{Model:} POYO+ \code{dim=128}, 6 self-attention layers, 16 latents per 125 ms, 2/8 cross/self heads, 5.2 M parameters.
\item \textbf{Optimiser:} LAMB, learning rate $3.125\times10^{-5}\times$ batch size, one-cycle cosine schedule,
  weight decay $10^{-4}$, gradient clipping at norm 1.
\item \textbf{Data:} 1-second windows, unit-dropout augmentation, cursor velocity z-scored with
  $\sigma = (5.66, 6.06)$ cm/s; annotated cursor-tracking artifacts ignored.
\item \textbf{Hardware:} Apple M1 Max GPU via PyTorch MPS; batch 64; \PretrainHours{} for \PretrainEpochs{} epochs.
\item \textbf{Known quirk:} right after each validation pass, PyTorch 2.2's MPS backend occasionally produces non-finite
  gradients in the encoder cross-attention (stale GPU buffers). Those steps are detected and skipped
  (\PretrainSkipped{} of \PretrainSteps{} during pretraining, $\approx$1\%), and the GPU cache is now released after
  validation.
\item \textbf{Monkey-T transfer:} 16 epochs, batch 32, LAMB at 10$\times$ the pretraining learning rate (new neuron
  embeddings start tiny and need bigger steps), transformer frozen for the first 6 epochs.
\end{itemize}

\subsection*{References}
\begin{itemize}
\item Azabou et al., \emph{A Unified, Scalable Framework for Neural Population Decoding} (POYO), NeurIPS 2023.
\item Azabou et al., \emph{Multi-session, multi-task neural decoding from distinct cell-types and brain regions} (POYO+), ICLR 2025.
\item Perich, Gallego \& Miller, \emph{A neural population mechanism for rapid learning}, Neuron 2018. Data: DANDI:000688.
\end{itemize}

\end{document}
